% Generated draft: requires mathematical and source review.
\documentclass[11pt]{article}
\usepackage[margin=25mm]{geometry}
\usepackage{fontspec}
\setmainfont{Latin Modern Roman}
\usepackage{amsmath,amssymb,mathtools,unicode-math}
\setmathfont{Latin Modern Math}
\usepackage{enumitem,needspace,xurl,hyperref}
\setlength{\parindent}{0pt}
\setlength{\parskip}{5pt}
\setlength{\emergencystretch}{3em}
\newcommand{\N}{\mathbb{N}}
\newcommand{\Z}{\mathbb{Z}}
\newcommand{\Q}{\mathbb{Q}}
\newcommand{\R}{\mathbb{R}}
\newcommand{\C}{\mathbb{C}}
\newcommand{\F}{\mathbb{F}}
\newcommand{\A}{\mathbb{A}}
\newcommand{\PP}{\mathbb{P}}
\newcommand{\E}{\mathbb{E}}
\newcommand{\eps}{\varepsilon}
\newcommand{\dd}{\,\mathrm{d}}
\begin{document}
\section*{UnsolvedMath ID: 6200031}
\textbf{Source problem code:} catalogue entry
\subsection{Definitions and mathematical statement}
Let \(G,H\) be finitely generated groups with word metrics. They are quasi-isometric if there is \(q:G\to H\) and constants \(L\ge1,C\ge0\) such that
\[
L^{-1}d_G(g,g')-C\le d_H(qg,qg')\le Ld_G(g,g')+C
\]
for all \(g,g'\), and every point of \(H\) is within \(C\) of \(q(G)\).

For a fixed \(n\ge2\), say \(G\) splits over \(\Z^n\) if it admits a nontrivial amalgamated-product or HNN decomposition with edge group isomorphic to \(\Z^n\). Equivalently, use a nontrivial Bass--Serre tree action without inversions and with one edge orbit whose edge stabilizers are isomorphic to \(\Z^n\). Papasoglu's first question is whether
\[
G\simeq_{\mathrm{QI}}H\quad\Longrightarrow\quad
\bigl(G\text{ splits over }\Z^n\iff H\text{ splits over }\Z^n\bigr),
\]
first for \(n=2\), and more generally for the intended ranks. This concerns existence of a splitting, not merely preservation of subgroups of a particular rank.

The source does not specify finite presentation, torsion, one-endedness, relative hypotheses, or whether edge groups are literally \(\Z^n\) or are intended only up to finite index/finite kernel. Exact and virtually free-abelian edge-group classes are different; quasi-isometry's insensitivity to finite extensions makes this distinction essential. No such replacement or restriction is silently made. The supplied discussion of virtually cyclic splittings for hyperbolic groups is a special-case comparison, not a hypothesis restricting the present higher-rank question to hyperbolic groups.

The second question asks for an analogous quasi-isometry invariance statement for JSJ decompositions. This requires specifying the allowed edge-group family, any relative subgroups, the JSJ equivalence convention, and what data a quasi-isometry should preserve, such as the decomposition tree or its deformation class. The extraction provides no particular JSJ object or precise matching assertion. The supplied research summary calls the general higher-rank problem open/triage without a verified resolution.
\subsection{Short English statement}
Does existence of a nontrivial splitting over \(\Z^2\), or higher-rank \(\Z^n\), pass between quasi-isometric groups? Formulate the corresponding JSJ question with an explicit edge-group class and equivalence convention.
\subsection{Sources}
\begin{itemize}
\item Local UnsolvedMath registry, permanent ID \(6200031\), exact source code \texttt{catalogue entry}, ``Boundaries of Groups and Kleinian Groups --- Problem 31.'' Proposer Panos Papasoglu, 2005.
\item Kapovich, \emph{Problems on Boundaries of Groups and Kleinian Groups} (2005), Problem 31, recorded PDF page 9; \url{https://www.math.ucdavis.edu/~kapovich/EPR/problems.pdf}. Recorded access 2026-07-29; pdf-text-normalized extraction.
\item Supplied OPEN-TRIAGE research summary mentions quasi-isometry invariance of virtually cyclic splittings for hyperbolic groups but gives no precise theorem or additional bibliographic citation.
\end{itemize}
\end{document}
